\section{Future Research Directions and Conclusion}
\label{sec:conclusion}

Although document intelligence has evolved from OCR-centered pipeline approaches to end-to-end systems capable of cross-page, multimodal understanding and retrieval-augmented generation, achieving stable deployment in real-world scenarios still faces critical challenges. These challenges lie not only in a model's ability to understand complex documents, but also in system-level evidence organization and reasoning mechanisms, and in risk-oriented evaluation paradigms for real-world deployment. Future research can be advanced along three levels.

\subsection{Verifiable and Structure-Aware Document Understanding}

Current document question answering systems can often produce seemingly plausible answers, yet these answers frequently lack explicit and verifiable evidence. For example, when asked about a key clause in a contract, a system may generate the correct conclusion but fail to indicate on which page or in which table cell the information is located. This phenomenon---``correct answers but opaque evidence''---limits adoption in high-stakes scenarios, because users cannot easily judge whether the model's reasoning process is trustworthy.

Therefore, a key research direction is to jointly model evidence localization and answer generation, so that while producing an answer, the model can explicitly identify the supporting document regions and establish a correspondence between the answer and its evidence, making the reasoning process verifiable rather than merely plausible.

Meanwhile, document understanding is shifting from pipeline methods that rely on explicit text extraction (e.g., OCR) to end-to-end vision-language models that operate directly on visual inputs. Such OCR-free methods can avoid information loss caused by recognition errors, but they also introduce new challenges: these models no longer have explicit textual and structural boundaries, and thus must automatically learn the document organization internally, such as the hierarchical relations between headings and body text, row and column alignments in tables, and semantic associations across fields.

Without such structure awareness, even if a model can recognize local content, it is difficult to perform consistent and reliable reasoning. Future research should develop structure-aware vision-language modeling methods, enabling models not only to ``see'' document content but also to understand how documents are organized, and to perform content recognition, structural inference, and evidence localization within a unified framework.

\subsection{Multimodal RAG and Agentic Systems for Complex Evidence Organization}

As document understanding tasks extend from single pages to long documents and even collections of multiple documents, the system's challenge is no longer merely to find the page that contains the answer, but to integrate information scattered across different locations into a consistent and reliable conclusion. For instance, an answer may partly come from textual descriptions, partly from tabular data, and partly from figures or supplementary notes on later pages.

Although current methods can retrieve relevant content, they often lack effective mechanisms to confirm whether these pieces of information refer to the same entity or are logically consistent, which can lead to incomplete or even contradictory reasoning results. Future research should develop finer-grained retrieval and representation methods that allow models to localize specific regions in documents, establish cross-page and cross-modal associations, and continuously verify consistency among different evidence during reasoning, rather than relying on one-shot retrieval results.

In addition, in real application scenarios, document understanding is typically not accomplished by a single model, but by the collaboration of multiple functional modules---for example, layout parsing to identify document structure, table parsing to extract structured data, retrieval modules to locate relevant content, and rule-based or computational modules to perform precise verification. Compared with a single model, an agentic system that integrates these tools can dynamically choose appropriate processing steps according to task requirements, thereby improving overall capability.

However, this also brings new research challenges, including how to ensure that the model can reliably select and invoke appropriate tools, how to prevent errors from accumulating throughout multi-step processing, and how to explicitly indicate that a reliable conclusion cannot be reached when evidence is insufficient or conflicting, rather than generating a superficially plausible but ungrounded answer.

\subsection{Risk-Oriented Data and Evaluation Paradigms for Real-World Deployment}

Most existing evaluations of document understanding systems are built on well-structured, high-quality benchmark data, such as well-typeset PDFs or form documents with fixed templates. In such environments, models can often achieve high accuracy. However, documents in real business scenarios are usually more complex: scans may be blurry, skewed, or occluded; documents from different sources can vary substantially in format and language; and key information may be scattered across long contexts or even multiple documents.

This gap between training and evaluation environments and real-world conditions leads to models that perform well on benchmarks but exhibit significant performance degradation in practice. Therefore, future evaluations should pay more attention to generalization under real distributions, including robustness to low-quality documents, adaptability to cross-domain documents, and computational efficiency and stability when handling long documents.

Moreover, existing evaluations typically focus only on whether the final answer is correct, but in practical applications an equally important question is whether the answer is grounded in verifiable evidence. For instance, even if a system occasionally produces the correct answer, users may still find it hard to trust the result if the system cannot explain where the evidence comes from.

Therefore, future evaluations should move beyond a single accuracy metric to multi-dimensional metrics that reflect system reliability, such as whether evidence is correctly localized, whether the reasoning process is internally consistent, whether the system can explicitly indicate uncertainty when it fails, and whether errors are diagnosable. Such risk-oriented evaluation paradigms will more accurately reflect a system's trustworthiness in real-world conditions, and will shift research goals from merely improving average accuracy to building document intelligence systems whose behavior is predictable, whose reasoning is verifiable, and whose behavior in real-world deployment is reliable.

\subsection{Conclusion}

This survey has provided a comprehensive review of document intelligence research, tracing the evolution from traditional OCR and layout analysis through deep learning innovations to contemporary multimodal foundation models. We have organized the field along the core pipeline of perception and reasoning, covering document layout analysis, OCR, table understanding, text-based retrieval-augmented generation, and vision-language models for document understanding.

Key trends emerge from this review: (1) the shift from modular pipelines to end-to-end unified models, (2) the increasing importance of structure-aware and verifiable reasoning, (3) the evolution from single-page to multi-page and long-document understanding, and (4) the growing emphasis on real-world robustness and deployment reliability. These trends point toward a future where document intelligence systems are not merely accurate on benchmarks, but trustworthy, interpretable, and reliable in practical applications.

The challenges ahead are substantial, requiring advances in model architectures, training paradigms, evaluation methodologies, and system integration. However, the rapid progress in foundation models, multimodal learning, and agentic systems provides strong foundations for addressing these challenges. We hope this survey serves as a useful resource for researchers and practitioners working to advance the field of document intelligence.
